\begin{table}[htbp]
\centering
\caption{Position of the dominant 600--850\,\wn{} band before and after promoting phosphorus, fluorine, sulfur and chlorine to a polarized triple-zeta basis.}
\label{tab:si-bandshift}
\small
\begin{tabular}{lrrrrrr}
\toprule
Agent & Measured & Baseline & Error & tight P & Error & Shift \\
\midrule
GF & 752 & 729 & -23 & 747 & -5 & +18 \\
GD & 752 & 748 & -3 & 762 & +10 & +14 \\
A-230 & 708 & 724 & +16 & 724 & +16 & -0 \\
A-232 & 708 & 658 & -50 & 663 & -45 & +5 \\
A-234 & 714 & 769 & +55 & 792 & +78 & +23 \\
HD & 733 & 743 & +10 & 746 & +12 & +3 \\
\bottomrule
\end{tabular}

\begin{flushleft}\footnotesize The shift is upward in 5 of the 6 molecules, by 3 to 23\,\wn{}, and downward by 0\,\wn{} in the remaining one, while the baseline errors scatter from $-$50 to $+$55\,\wn{}. A correction of fixed sign can only help the subset whose error is negative and of comparable size. GF satisfies that condition; GD, which differs only in the ester substituent and was already within 3\,\wn{}, does not.\end{flushleft}
\end{table}
